\documentclass[12pt]{article}
\usepackage[utf8]{inputenc}
\usepackage[T1]{fontenc}
\usepackage{hyperref}
\title{Fang Language Resources for K-12 Curriculum Development}
\author{Compiled by Claude Sonnet 5.5 (claude-sonnet-5-5) for Andr\'as Kornai}
\date{}
\begin{document}
\maketitle

\section{Introduction}
Fang of Cameroon (ISO 639-3 \texttt{fak}) is distinct from the much larger Fang language (\texttt{fan}) of Equatorial Guinea and Gabon; sources often confuse the two, and resources for \texttt{fan} must not be assumed to apply. Speaker figures vary widely between sources (5,900 in one Joshua Project profile, 111,000 in a UNESCO listing). Search found no dictionary, grammar, primer or Bible translation specific to Fang (Cameroon).

\section{Primary Educational Resources}
None found.

\section{Linguistic and Reference Materials}
\begin{itemize}
\item Wikipedia on Fang (Cameroon): \url{https://en.wikipedia.org/wiki/Fang_language_(Cameroon)}
\item UNESCO World Atlas of Languages entry: \url{https://ar.wal.unesco.org/en/countries/cameroon/languages/fang-cameroon}
\item PHOIBLE phoneme inventory for \texttt{fak}: \url{https://phoible.org/languages/fak}
\end{itemize}

\section{Religious and Community Resources}
\begin{itemize}
\item Joshua Project, Benjube Fenga Fang in Cameroon (Bible translation status: needed): \url{https://www.joshuaproject.net/people_groups/19661/CM}
\item Joshua Project, Ntum Fang in Cameroon (states Scripture has not been translated into this Fang): \url{https://joshuaproject.net/index.php/people_groups/11187/CM}
\end{itemize}

\section{Digital and Audio Resources}
None specific to Fang (Cameroon) were identified.

\section{Curriculum Development Context}
Documentation is minimal. Clarifying which speech varieties are meant, and their relation to the neighbouring Beti languages, is a prerequisite to curriculum work.

\end{document}